\documentclass[11pt,a4paper]{article}
\usepackage[utf8]{inputenc}
\usepackage[T1]{fontenc}
\usepackage[french]{babel}
\usepackage{geometry}
\usepackage{titlesec}
\usepackage{parskip}
\usepackage{enumitem}
\usepackage{listings}
\usepackage{xcolor}
\usepackage{booktabs}
\usepackage{fancyhdr}
\usepackage{amsmath}
\usepackage{amssymb}

\geometry{margin=2.2cm}

% Colors
\definecolor{nkblue}{RGB}{25, 60, 120}
\definecolor{nkgray}{RGB}{245, 245, 248}
\definecolor{codegray}{RGB}{240, 240, 240}
\definecolor{codegreen}{RGB}{0, 120, 50}
\definecolor{codepurple}{RGB}{120, 0, 120}

% Section formatting
\titleformat{\section}{\large\bfseries\color{nkblue}}{}{0em}{\thesection\quad}[\vspace{-4pt}\rule{\linewidth}{0.6pt}]
\titleformat{\subsection}{\normalsize\bfseries}{}{0em}{\thesubsection\quad}

% Code style
\lstset{
  backgroundcolor=\color{codegray},
  basicstyle=\ttfamily\footnotesize,
  keywordstyle=\color{nkblue}\bfseries,
  commentstyle=\color{codegreen},
  stringstyle=\color{codepurple},
  breaklines=true,
  frame=single,
  rulecolor=\color{gray!40},
  framesep=4pt,
  language=C++
}

% Header / Footer
\pagestyle{fancy}
\fancyhf{}
\fancyhead[L]{\small\textbf{Réalité Augmentée — NkMath from scratch}}
\fancyhead[R]{\small Semestre 1 · Semaines 1–4}
\fancyfoot[C]{\small\thepage}
\renewcommand{\headrulewidth}{0.4pt}

% --------------------------------------------------
\begin{document}

\begin{titlepage}
  \centering
  \vspace*{2cm}
  {\Huge\bfseries\color{nkblue} Réalité Augmentée\\[6pt] \textit{from scratch}\par}
  \vspace{0.5cm}
  \rule{0.6\linewidth}{1.2pt}
  \vspace{0.5cm}

  {\large Rapport de synthèse — Semestres 1, Semaines 1 à 4\par}
  \vspace{0.3cm}
  {\normalsize Modules 01 à 04 · Bibliothèque mathématique \texttt{NkMath}\par}
  \vspace{1cm}
  {\small\textit{Règle absolue : Zéro bibliothèque tierce — C++17 + STL uniquement}\par}
  \vspace{2cm}

  \begin{tabular}{ll}
    \textbf{Auteur du cours} & Teuguia Tadjuidje Rodolf Sédéris · Rihen \\
    \textbf{Durée totale}    & 3 Semestres · 162h · 21 Modules \\
    \textbf{Modules couverts} & M01 – M04 (Semaines 1–4) \\
  \end{tabular}
  \vfill
\end{titlepage}

% --------------------------------------------------
\section*{Introduction}
Ce rapport résume les quatre premières semaines du cours \textit{Réalité Augmentée from scratch}.
L'objectif de ce semestre est de construire la bibliothèque mathématique \texttt{NkMath} intégralement
en C++17, sans aucune dépendance externe. Chaque module pose une brique indispensable au
pipeline AR complet : arithmétique flottante fiable, algèbre vectorielle, transformations matricielles
et représentation des rotations par quaternions.

\tableofcontents
\newpage

% ==================================================
\section{Semaine 1 — Module 01 : IEEE 754 \& Arithmétique Flottante}
\textit{Durée : 6 heures}

\subsection{Pourquoi ce module est fondamental}
Tout le pipeline AR manipule des flottants : coordonnées 3D, matrices de projection, résidus de
calibration, valeurs de pixels. Une erreur d'arrondi invisible peut faire diverger des algorithmes
itératifs (Levenberg-Marquardt), produire une homographie incorrecte, ou décaler une pose de
plusieurs centimètres. Ce module est la fondation de tout code numérique fiable.

\subsection{Représentation binaire IEEE 754}
Un \texttt{float} 32 bits est découpé en trois champs :

\begin{center}
\begin{tabular}{lll}
  \toprule
  \textbf{Champ} & \textbf{Largeur} & \textbf{Rôle} \\
  \midrule
  Signe    & 1 bit   & 0 = positif, 1 = négatif \\
  Exposant & 8 bits  & Valeur biaisée (biais = 127) \\
  Mantisse & 23 bits & Fraction avec bit implicite \texttt{1.} \\
  \bottomrule
\end{tabular}
\end{center}

\[
  \text{Valeur} = (-1)^{\text{signe}} \times 2^{\text{exposant} - 127} \times 1.\text{mantisse}
\]

Un \texttt{double} 64 bits porte 11 bits d'exposant et 52 bits de mantisse, soit environ 15–16 chiffres
décimaux significatifs (contre 7–8 pour le \texttt{float}).

En C++17, la seule façon légale d'accéder aux bits d'un flottant est via \texttt{std::memcpy} vers un
entier, afin d'éviter tout comportement indéfini lié au \textit{strict aliasing}.

\subsection{Valeurs spéciales}
\begin{itemize}[noitemsep]
  \item \textbf{$\pm$Inf} : exposant \texttt{0xFF}, mantisse nulle (ex. \texttt{1.0f/0.0f}).
  \item \textbf{NaN} : exposant \texttt{0xFF}, mantisse non nulle ; \textit{contagieux} — toute opération
        avec NaN produit NaN. La seule comparaison valide est \texttt{std::isnan(x)}, car \texttt{NaN != NaN} est toujours vrai.
  \item \textbf{Subnormaux} : exposant \texttt{0x00}, mantisse non nulle — représentation des très
        petits nombres au prix d'une perte de précision.
  \item \textbf{$\pm$0} : deux patterns de bits distincts, mais mathématiquement égaux.
\end{itemize}

\subsection{Epsilon machine et ULP}
\textbf{Epsilon machine} ($\varepsilon$) : plus petite valeur telle que $1.0 + \varepsilon \neq 1.0$.
Pour \texttt{float}, $\varepsilon \approx 1.19\times10^{-7}$ ; pour \texttt{double},
$\varepsilon \approx 2.22\times10^{-16}$.

\textbf{ULP} (\textit{Unit in the Last Place}) : espacement entre deux flottants consécutifs au
voisinage d'une valeur donnée. L'ULP \textit{varie} avec la magnitude : près de 1000, l'ULP est
environ $1.22\times10^{-4}$, soit $10^3$ fois plus grand que près de 1.

\subsection{Cancellation catastrophique}
La soustraction de deux flottants presque égaux détruit presque tous les bits significatifs du
résultat. C'est l'erreur la plus dangereuse du pipeline car elle est silencieuse (pas de NaN,
pas d'Inf). Exemple : calculer la variance par la formule naïve
$\frac{\sum x_i^2 - (\sum x_i)^2/n}{n-1}$
peut produire une valeur \textit{négative} à cause de l'annulation.

\textbf{Remède — Algorithme de Welford} : mise à jour incrémentale de la moyenne et de la
variance qui n'effectue jamais de soustraction de grandes valeurs proches.

\subsection{Sommation de Kahan}
L'algorithme de Kahan compense les erreurs d'arrondi lors de la sommation de nombreux termes
grâce à une variable de compensation qui capture les bits perdus à chaque addition. Sur un million
de \texttt{0.1f}, \texttt{std::accumulate} donne une erreur de l'ordre de 11, tandis que Kahan
donne un résultat correct à 6 chiffres significatifs.

\subsection{Stratégie float/double dans NkMath}
\begin{itemize}[noitemsep]
  \item Calculs internes (matrices, SVD) : \textbf{double} pour la précision de calibration.
  \item Upload GPU (\texttt{glVertexAttribPointer}, \texttt{glUniformMatrix4fv}) : \textbf{float},
        car GLSL est optimisé pour le 32 bits.
  \item Types exportés : \texttt{Vec3d} (double) et \texttt{Vec3f} (float) selon le contexte.
\end{itemize}

\subsection{Travaux pratiques}
Les TP de cette semaine demandent : (1) implémenter \texttt{inspectFloat} pour visualiser les trois
champs IEEE 754 de valeurs cibles ; (2) démontrer empiriquement les pertes de précision avec
Kahan et Welford ; (3) créer le fichier \texttt{NkMath/Float.h} contenant toutes les fonctions
utilitaires (\texttt{isFiniteValid}, \texttt{nearlyZero}, \texttt{approxEq}, \texttt{kahanSum}) avec
au moins 30 tests unitaires.

% ==================================================
\section{Semaine 2 — Module 02 : Vecteurs — Vec2d, Vec3d, Vec4d}
\textit{Durée : 9 heures}

\subsection{Rôle des vecteurs dans le pipeline AR}
Les vecteurs sont la structure de données la plus utilisée du pipeline. Chaque position dans
l'espace est un \texttt{Vec3d}, chaque pixel est un \texttt{Vec2d}, chaque coordonnée homogène
est un \texttt{Vec4d}. Un tableau de \texttt{Vec3d} peut être passé directement à
\texttt{glVertexAttribPointer} \textit{sans copie} si le layout mémoire est correct, d'où l'usage
de \texttt{static\_assert} sur \texttt{sizeof} et \texttt{offsetof}.

\subsection{Vec2d}
Structure avec deux \texttt{double} (x, y) continus en mémoire. Fournit :
\begin{itemize}[noitemsep]
  \item Opérateurs arithmétiques : \texttt{+}, \texttt{-}, \texttt{*}, \texttt{/}, composés, unaire.
  \item Accès par index via \texttt{operator[]} (utilise le fait que les membres sont contigus).
  \item \texttt{Norm()}, \texttt{Normalized()}, \texttt{IsNormalized()}.
  \item \textbf{Dot product} : $a \cdot b = a_x b_x + a_y b_y = |a||b|\cos\theta$.
  \item \textbf{Cross 2D} (scalaire) : $a_x b_y - a_y b_x$ — signe indique si $b$ est à gauche
        ou à droite de $a$.
  \item \textbf{Lerp} linéaire.
\end{itemize}
Garantie : \texttt{static\_assert(sizeof(Vec2d)==16)}.

\subsection{Dot product — interprétation géométrique}
\[
  a \cdot b = |a||b|\cos\theta
  \begin{cases}
    = 0  & \Rightarrow \text{perpendiculaires} \\
    > 0  & \Rightarrow \text{angle aigu} \\
    < 0  & \Rightarrow \text{angle obtus}
  \end{cases}
\]
Usages dans le pipeline : test face avant/arrière, éclairage de Phong, projection d'un point sur
une droite, test d'appartenance à un plan.

\subsection{Vec3d et cross product}
Le \textbf{produit vectoriel} $a \times b$ produit un vecteur perpendiculaire à $a$ et $b$, dont la
norme est $|a||b|\sin\theta$ (aire du parallélogramme). Il respecte la règle de la main droite et
est \textit{non-commutatif} : $a \times b = -(b \times a)$.

Formule explicite :
\[
  (a \times b) = \begin{pmatrix} a_y b_z - a_z b_y \\ a_z b_x - a_x b_z \\ a_x b_y - a_y b_x \end{pmatrix}
\]

\textbf{Projection et rejection} : décomposer un vecteur $a$ en composante parallèle et
perpendiculaire à $b$ :
\[
  \text{proj}(a,b) = b \,\frac{a \cdot b}{|b|^2}, \qquad \text{reject}(a,b) = a - \text{proj}(a,b)
\]

\subsection{Gram-Schmidt — base orthonormale}
À partir de trois vecteurs quelconques (non coplanaires), l'algorithme de Gram-Schmidt construit
une base orthonormale $\{u, v, w\}$ en trois étapes : normaliser le premier vecteur, ôter sa
composante du second puis normaliser, ôter les deux composantes du troisième puis normaliser.
Utilisé pour construire la matrice \textit{LookAt} et la décomposition polaire (Module 05).

\subsection{Vec4d — coordonnées homogènes}
Les coordonnées homogènes unifient points et directions :
\begin{itemize}[noitemsep]
  \item $w = 1$ : \textbf{point} dans l'espace (la translation s'applique).
  \item $w = 0$ : \textbf{direction} (la translation ne s'applique pas).
\end{itemize}
La déhomogénéisation se fait par \texttt{ToVec3()} : diviser $(x, y, z)$ par $w$.

\subsection{Travaux pratiques}
Les TP demandent : (1) implémenter \texttt{Vec2d} complet avec 20 tests unitaires sur le dot
product, le cross 2D, la normalisation et l'accès par index ; (2) implémenter \texttt{Vec3d} avec
cross product et Gram-Schmidt, testés sur des triplets aléatoires ; (3) projeter les 8 coins d'un
cube unitaire en 2D via une projection perspective simple et tracer les 12 arêtes dans une image
PPM 512×512.

% ==================================================
\section{Semaine 3 — Module 03 : Matrices — Mat3d, Mat4d}
\textit{Durée : 9 heures}

\subsection{Convention column-major — critique pour OpenGL}
OpenGL stocke les matrices en \textbf{column-major} : dans un tableau linéaire, les éléments de la
première colonne viennent en premier. L'accès C++ est \texttt{data[col*4 + row]}. Lors de
l'upload GPU, \texttt{glUniformMatrix4fv} avec \texttt{GL\_FALSE} suppose déjà du column-major.
Une transposition accidentelle produit des transformations complètement incorrectes.

\subsection{Mat4d — structure et produit}
\begin{itemize}[noitemsep]
  \item Stockage : tableau \texttt{double data[16]} column-major.
  \item \texttt{operator()(row, col)} : accès intuitif en notation mathématique.
  \item Produit matriciel $4\times4$ : 64 multiplications et 48 additions.
  \item \texttt{operator*(Vec4d)} : transformation d'un vecteur homogène.
  \item \texttt{Transposed()}, \texttt{ToFloat(float[16])} pour l'upload GPU.
\end{itemize}

\subsection{Inverse par Gauss-Jordan avec pivot partiel}
L'inversion opère sur la matrice augmentée $[M \mid I]$ :
\begin{enumerate}[noitemsep]
  \item Pour chaque colonne, trouver la ligne avec le plus grand coefficient en valeur absolue
        (\textbf{pivot partiel}) et l'amener en position de pivot — réduit les erreurs numériques.
  \item Normaliser la ligne pivot.
  \item Éliminer la colonne courante dans toutes les autres lignes.
  \item Extraire la partie droite $I \to M^{-1}$.
\end{enumerate}
La fonction retourne \texttt{false} si la matrice est singulière.

\subsection{Matrice de rotation — formule de Rodrigues}
Rotation d'angle $\theta$ autour d'un axe unitaire $\hat{n} = (n_x, n_y, n_z)$ :
\[
  R = I\cos\theta + (1-\cos\theta)\,\hat{n}\hat{n}^T + \sin\theta\,[\hat{n}]_\times
\]
En pratique, chacun des 9 coefficients est calculé explicitement à partir de $c = \cos\theta$,
$s = \sin\theta$, $t = 1-c$.

\subsection{Matrice View — LookAt}
Construit la vue caméra à partir de trois vecteurs de l'espace monde :
\begin{enumerate}[noitemsep]
  \item $\hat{f} = \text{normalize}(\text{target} - \text{eye})$ — axe forward.
  \item $\hat{r} = \text{normalize}(\hat{f} \times \text{worldUp})$ — axe right.
  \item $\hat{u} = \hat{r} \times \hat{f}$ — axe up corrigé.
\end{enumerate}
La matrice View combine la rotation $[\hat{r}, \hat{u}, -\hat{f}]$ et la translation
$(-\hat{r}\cdot\text{eye},\; -\hat{u}\cdot\text{eye},\; \hat{f}\cdot\text{eye})$.

\subsection{TRS — Translation, Rotation, Scale}
\[
  M_{\text{TRS}} = T \times R \times S
\]
L'ordre est critique : d'abord l'échelle (locale), puis la rotation, puis la translation monde.
Inverser cet ordre produit des transformations incorrectes (ex. une mise à l'échelle après rotation
déforme l'objet dans l'espace monde).

\subsection{Travaux pratiques}
Les TP demandent : (1) implémenter \texttt{Mat4d} avec produit et inversion, vérifier
$M \times M^{-1} = I$ ; (2) construire un rasteriseur logiciel dans une image PPM 512×512
affichant les arêtes d'un cube animé via LookAt et projection perspective ; (3) implémenter TRS
complet et vérifier la décomposition inverse sur 20 triplets aléatoires.

% ==================================================
\section{Semaine 4 — Module 04 : Quaternions}
\textit{Durée : 6 heures}

\subsection{Problème du gimbal lock}
Les angles d'Euler (yaw, pitch, roll) souffrent du \textbf{gimbal lock} : lorsque le pitch
atteint $\pm 90°$, les axes yaw et roll deviennent coplanaires et le système perd un degré de
liberté. De plus, l'interpolation LERP entre deux orientations Euler ne suit pas le chemin le plus
court sur la sphère des rotations, ce qui produit des mouvements non-naturels.

\subsection{Définition et structure}
Un quaternion $q = w + xi + yj + zk$ avec $i^2 = j^2 = k^2 = ijk = -1$.
Un \textbf{quaternion unitaire} ($\|q\| = 1$) représente une rotation de $2\arccos(w)$ radians
autour de l'axe $(x, y, z)$.

Opérations fondamentales :
\begin{itemize}[noitemsep]
  \item \textbf{Conjugué} : $q^* = (w, -x, -y, -z)$ — inverse la rotation.
  \item \textbf{Produit de Hamilton} : composition de rotations (non-commutatif).
  \item Rotation d'un vecteur : $v' = q \otimes (0, v) \otimes q^*$, optimisée sans construire deux
        quaternions intermédiaires.
\end{itemize}

Conversion axe-angle $\to$ quaternion :
\[
  q = \left(\cos\frac{\theta}{2},\; \hat{n}\sin\frac{\theta}{2}\right)
\]

\subsection{SLERP — interpolation sphérique}
Le \textbf{SLERP} (\textit{Spherical Linear intERPolation}) suit le chemin le plus court sur la
sphère des quaternions unitaires avec une vitesse angulaire \textit{constante} :
\[
  \text{Slerp}(a, b, t) = \frac{\sin((1-t)\Omega)}{\sin\Omega}\,a + \frac{\sin(t\Omega)}{\sin\Omega}\,b, \qquad \Omega = \arccos(a \cdot b)
\]
Lorsque $\Omega$ est très petit ($\cos\Omega > 0.9999$), on bascule vers un LERP normalisé pour
éviter la division par zéro. Si $a \cdot b < 0$, on inverse $b$ pour prendre le chemin court.

\subsection{Méthode de Shepperd — conversion Quat $\leftrightarrow$ Mat3}
La conversion quaternion $\to$ matrice $3\times3$ est directe (9 produits). La conversion inverse
(\textbf{méthode de Shepperd}) est numériquement stable car elle choisit, selon la trace de $R$,
la composante dominante du quaternion pour éviter les divisions par de petites valeurs :
\begin{itemize}[noitemsep]
  \item \texttt{trace > 0} : calculer $w$ en priorité.
  \item \texttt{R(0,0)} dominant : calculer $x$ en priorité.
  \item Idem pour les cas $y$ et $z$ dominants.
\end{itemize}

\subsection{Travaux pratiques}
Les TP demandent : (1) valider les quaternions (rotation, aller-retour Quat/Mat3, inverse) ;
(2) animer en 60 frames une rotation entre deux orientations et comparer visuellement SLERP
(vitesse angulaire uniforme) et LERP (vitesse non-uniforme) dans le rasteriseur logiciel du Module 03.

% ==================================================
\section*{Synthèse}
\begin{center}
\begin{tabular}{llll}
  \toprule
  \textbf{Semaine} & \textbf{Module} & \textbf{Durée} & \textbf{Apport clé} \\
  \midrule
  1 & IEEE 754 \& flottants    & 6h & Précision numérique, Kahan, Welford \\
  2 & Vec2d, Vec3d, Vec4d      & 9h & Algèbre vectorielle, dot, cross, Gram-Schmidt \\
  3 & Mat3d, Mat4d             & 9h & Transformations, Rodrigues, LookAt, TRS \\
  4 & Quaternions              & 6h & Rotations stables, SLERP, Shepperd \\
  \midrule
  \multicolumn{2}{l}{\textbf{Total S1 (4 semaines)}} & \textbf{30h} & Bibliothèque \texttt{NkMath} — fondations \\
  \bottomrule
\end{tabular}
\end{center}

Ces quatre modules forment le socle mathématique sur lequel reposent tous les algorithmes du
pipeline AR : SVD (M05), détection de features (M07–08), calibration (M10–11), estimation de
pose PnP (M15) et rendu OpenGL (M18–20).

\end{document}
